\chapter{Volet Business Intelligence et Prévision}

\section*{Introduction}
Au-delà de la prédiction ponctuelle d'une empreinte environnementale, GREENLY vise à donner à l'utilisateur une vision analytique et évolutive de son impact. Ce quatrième chapitre présente le second grand volet du projet~: le module de \textbf{Business Intelligence} (BI), qui transforme l'historique des analyses individuelles en indicateurs, tableaux de bord et prévisions exploitables. Nous y détaillons l'architecture décisionnelle retenue, les indicateurs clés de performance calculés, les tableaux de bord et visualisations développés, le module de prévision statistique, ainsi que le système de recommandations personnalisées.

\section{Architecture BI}

Contrairement à une architecture décisionnelle classique reposant sur un entrepôt de données (\textit{data warehouse}) séparé et des traitements ETL périodiques, GREENLY adopte une architecture BI \textbf{légère et temps réel}, adaptée à l'échelle d'un historique individuel~:

\begin{itemize}[label=\textbullet,font=\normalsize]
    \item la table \texttt{analyses} de la base Supabase joue le rôle de \textbf{table de faits}~: chaque ligne représente un événement de mesure (une analyse), avec ses dimensions (catégorie, date) et ses mesures (\texttt{co2e}, \texttt{green\_score}, contenu JSON \texttt{outputs})~;
    \item l'\textbf{agrégation} (regroupement mensuel, annuel, par catégorie) est calculée \textbf{côté client}, en React, via des fonctions mémoïsées (\texttt{useMemo}) qui recalculent les indicateurs à chaque changement de filtre (période, catégorie) sans requête serveur supplémentaire~;
    \item le \textbf{calcul prédictif} (prévision de tendance) est délégué au micro-service ML via l'endpoint \texttt{POST~/api/forecast}, qui applique une régression statistique sur les points d'historique transmis par le client.
\end{itemize}

Ce choix d'architecture évite la complexité d'un entrepôt de données dédié, injustifiée à l'échelle des volumes traités (historique d'un utilisateur individuel), tout en conservant la réactivité nécessaire à une expérience de tableau de bord interactif.

\section{KPIs}

Le tableau~\ref{tab:kpis} présente les indicateurs clés de performance (\textit{Key Performance Indicators}) calculés et affichés par la plateforme.

\begin{table}[H]
\centering
\caption{Indicateurs clés de performance (KPI) de GREENLY}
\label{tab:kpis}
\begin{tabular}{|p{3.4cm}|p{6.4cm}|p{3.5cm}|}
\hline
\rowcolor{isiBlue!15}
\textbf{KPI} & \textbf{Définition} & \textbf{Mode de calcul} \\
\hline
Green Score & Score de 0 à 100 situant une analyse par rapport à la distribution des valeurs observées à l'entraînement. & $100 - \text{percentile}$ (voir chapitre~3) \\ \hline
CO2e total & Somme des émissions de CO2 équivalent sur la période et les catégories sélectionnées. & $\sum \texttt{co2e}$ \\ \hline
Consommation d'eau & Somme des litres d'eau calculés sur la période sélectionnée. & $\sum \texttt{waterLiters}$ \\ \hline
Consommation d'énergie & Somme des kWh calculés (énergie domestique + électronique). & $\sum \texttt{energyKwh}$ \\ \hline
Score moyen & Moyenne des \textit{Green Scores} sur la période/catégorie sélectionnée. & $\overline{\texttt{green\_score}}$ \\ \hline
Arbres nécessaires & Nombre d'arbres qu'il faudrait planter pour compenser le CO2e total annuel (1 arbre~$\approx$~21~kg CO2/an). & $\lceil \text{CO2e total} / 21 \rceil$ \\ \hline
Confiance du modèle & Fiabilité du modèle de prédiction de la catégorie active, exprimée en \%. & $R^2_{\text{modèle}} \times 100$ \\ \hline
Erreur de prévision (MAE) & Erreur absolue moyenne du modèle sur le jeu de test, affichée pour contextualiser la confiance accordée à une prédiction. & Voir tableau~\ref{tab:resultats-globaux} \\ \hline
\end{tabular}
\end{table}

\section{Tableaux de bord}

GREENLY propose trois tableaux de bord complémentaires, correspondant à trois niveaux de lecture de l'impact environnemental de l'utilisateur~:

\begin{enumerate}
    \item \textbf{Dashboard (vue d'ensemble)}~: page d'accueil de l'espace utilisateur, présentant les KPI principaux, une jauge de score de durabilité, la tendance mensuelle du CO2e, la répartition des analyses par catégorie, ainsi que des \textit{insights} textuels générés automatiquement (comparaison au mois précédent, catégorie la plus suivie).
    \item \textbf{Tableau de bord BI}~: vue analytique avancée avec filtres par période (3, 6, 12 mois, année en cours) et par catégorie, comparaison annuelle, radar de score par catégorie face à un objectif cible, et une synthèse exécutive générée à partir des agrégats réels de l'utilisateur.
    \item \textbf{AI Analytics}~: vue centrée sur les modèles eux-mêmes~: métriques réelles de performance par catégorie, comparaison des 5 algorithmes, importance des variables, prévision de tendance et recommandations dérivées.
\end{enumerate}

\section{Visualisations}

L'ensemble des visualisations est développé avec la bibliothèque \textbf{Recharts} (React), permettant des graphiques interactifs (infobulles, légendes, mise en évidence au survol) directement intégrés à l'interface. Six familles de représentations graphiques sont utilisées~: graphiques en aires (tendance temporelle), histogrammes (comparaison catégorielle ou annuelle), diagrammes circulaires (répartition), diagrammes radar (score multi-catégoriel face à un objectif), courbes (évolution du score) et graphiques composés (prévision avec intervalle de confiance).

La figure~\ref{fig:bi-r2} présente une donnée réelle du projet~: la précision (R\textsuperscript{2}) du modèle finalement retenu pour chacune des six catégories, telle qu'exposée par l'API \texttt{GET~/api/models/metrics} et affichée sur la page AI Analytics.

\begin{figure}[H]
\centering
\includegraphics[width=0.82\textwidth]{bi-r2-par-categorie}
\caption{Précision réelle (R\textsuperscript{2}) du modèle retenu par catégorie (donnée réelle \texttt{metrics.json})}
\label{fig:bi-r2}
\end{figure}

Les figures~\ref{fig:bi-tendance} et \ref{fig:bi-repartition} illustrent, à partir de données d'exemple (l'application ne disposant pas encore d'historique utilisateur en production au moment de la rédaction de ce rapport), le rendu des deux graphiques principaux du tableau de bord BI~: la tendance mensuelle multi-indicateurs et la répartition des analyses par catégorie.

\begin{figure}[H]
\centering
\includegraphics[width=0.85\textwidth]{bi-tendance-mensuelle}
\caption{Exemple de rendu -- tendance mensuelle CO2e / énergie (données d'illustration)}
\label{fig:bi-tendance}
\end{figure}

\begin{figure}[H]
\centering
\includegraphics[width=0.6\textwidth]{bi-repartition-categories}
\caption{Exemple de rendu -- répartition des analyses par catégorie (données d'illustration)}
\label{fig:bi-repartition}
\end{figure}

La figure~\ref{fig:bi-radar} illustre le graphique radar comparant, pour chaque catégorie, le score moyen de l'utilisateur à un objectif fixé à 85/100 -- une représentation particulièrement efficace pour identifier en un coup d'œil les catégories prioritaires d'amélioration.

\begin{figure}[H]
\centering
\includegraphics[width=0.55\textwidth]{bi-radar-score}
\caption{Exemple de rendu -- Green Score par catégorie face à l'objectif (données d'illustration)}
\label{fig:bi-radar}
\end{figure}

\section{Prévisions}

Le module de prévision (\texttt{app/forecasting.py}) projette la tendance future d'un utilisateur à partir de son historique réel de \textit{Green Scores} ou de valeurs d'empreinte, catégorie par catégorie. Volontairement, ce module \textbf{n'utilise pas} un modèle de \textit{Machine Learning} complexe ni de saisonnalité simulée~: le service ML n'a pas accès à l'historique complet de l'utilisateur (celui-ci résidant côté Supabase/frontend), et une régression linéaire simple sur les points réellement transmis par le client est jugée plus honnête qu'un modèle sophistiqué appliqué à un historique potentiellement très court.

La méthode ajuste une régression linéaire $y = a x + b$ par la méthode des moindres carrés sur les $n$ points d'historique, puis extrapole sur les périodes futures demandées, avec un intervalle de confiance qui \textbf{s'élargit avec l'horizon de projection}~:

\begin{lstlisting}[language=Python]
slope, intercept = np.polyfit(x, values, 1)
fitted = slope * x + intercept
residual_std = np.std(values - fitted)
...
widen = residual_std * (1 + 0.15 * i)   # incertitude croissante avec l'horizon
predicted = slope * future_x + intercept
\end{lstlisting}

Le niveau de confiance accordé à la prévision (\texttt{modelConfidence}) n'est pas une constante arbitraire~: il combine la qualité de l'ajustement (un pseudo-R\textsuperscript{2} calculé sur les résidus) et la taille de l'échantillon disponible (une prévision fondée sur 2 points est nécessairement moins fiable qu'une prévision fondée sur 12 points)~:

\begin{lstlisting}[language=Python]
r2 = 1 - (np.var(residuals) / np.var(values))
sample_factor = min(1.0, n / 12)
model_confidence = 0.5 * max(0.0, r2) + 0.5 * sample_factor
\end{lstlisting}

La figure~\ref{fig:bi-forecast} présente une exécution \textbf{réelle} de cet algorithme (code non modifié, exécuté tel quel), appliquée à un historique d'exemple de huit mois de CO2e mensuel. La tendance détectée est \textit{décroissante} ($-14{,}4\,\%$ projetés sur les 4 prochains mois), illustrant une situation où l'empreinte de l'utilisateur s'améliore dans le temps.

\begin{figure}[H]
\centering
\includegraphics[width=0.85\textwidth]{bi-forecast}
\caption{Prévision de tendance -- exécution réelle de l'algorithme sur un historique d'exemple}
\label{fig:bi-forecast}
\end{figure}

\section{Analyses}

Au-delà des graphiques, le tableau de bord BI génère une \textbf{synthèse analytique automatique} (« AI Executive Summary ») directement calculée à partir des agrégats réels de l'utilisateur, sans recours à un service de génération de texte externe~: la meilleure et la moins bonne catégorie sont déterminées par comparaison des scores moyens, et la tendance globale est établie en comparant la moyenne des scores sur la première et la seconde moitié de la période sélectionnée~:

\begin{lstlisting}[language=Python]
sorted_scores = sorted(scoreByCategory, key=lambda s: -s.current)
best, worst = sorted_scores[0], sorted_scores[-1]
improving = secondHalfAvgScore >= firstHalfAvgScore
\end{lstlisting}

Cette synthèse restitue par exemple des messages tels que \textit{« Votre score moyen s'est amélioré de 6 points sur cette période »} ou \textit{« Transport est votre catégorie la plus performante, avec une moyenne de 82/100 »}, personnalisés à chaque utilisateur et recalculés à chaque changement de filtre.

\section{Recommandations}

Le système de recommandations personnalisées constitue l'un des apports les plus directement liés au volet Intelligence Artificielle du chapitre précédent~: plutôt que d'afficher des conseils génériques identiques pour tous les utilisateurs (limite identifiée chez les solutions concurrentes, chapitre~1), GREENLY exploite l'\textbf{importance réelle des variables} (\textit{feature importance}) calculée par chaque modèle entraîné pour cibler la recommandation la plus pertinente pour la situation spécifique de l'utilisateur~:

\begin{lstlisting}[language=Python]
tips = sorted(feature_importance.items(), key=lambda kv: -kv[1])[:3]
# ex. pour la categorie "carbon" :
# distance (43%) -> "La distance domine votre empreinte -- grouper vos
#                     trajets ou le télétravail change beaucoup la donne."
# fuelType (27%) -> "Le type de carburant est un facteur déterminant --
#                     passer à l'hybride ou à l'électrique réduirait..."
\end{lstlisting}

Chaque recommandation est ainsi accompagnée du poids réel (en \%) que la variable correspondante représente dans les prédictions du modèle pour la catégorie concernée, renforçant la transparence et la crédibilité du conseil donné -- un utilisateur voit explicitement \textit{pourquoi} telle variable lui est signalée, plutôt que de recevoir un conseil non justifié.

\section*{Conclusion}
Ce chapitre a présenté le volet Business Intelligence et prévision de GREENLY~: une architecture décisionnelle légère fondée sur l'agrégation côté client, des indicateurs clés de performance directement dérivés des sorties des modèles de Machine Learning, trois tableaux de bord complémentaires, un module de prévision statistique honnête quant à son incertitude, ainsi qu'un système de recommandations personnalisées fondé sur l'importance réelle des variables. Combiné au volet Intelligence Artificielle détaillé au chapitre~3, ce module de Business Intelligence positionne GREENLY comme une plateforme de suivi environnemental complète, dépassant le simple calculateur ponctuel pour offrir un véritable accompagnement dans la durée. Le chapitre suivant conclut ce rapport en dressant le bilan du travail réalisé et en proposant des perspectives d'évolution.
